\chapter{Envelope systems over Hermitian perturbations}
\label{chap:hermitian-envelope-systems}

Enlarging a background allows more plurisubharmonic potentials, but multiplier-ideal projection retains the singularity data of each test-curve slice. Compatible projections over all positive Hermitian perturbations form an envelope system. As the perturbation shrinks, their normalized slices decrease to a curve in the original background. This limit recovers the system precisely when no multiplier ideals are lost. A ruled surface will show how an additional divisorial pole can obstruct recovery even in a K\"ahler background.

We use curves that are constant on a fixed left half-line, so their initial potential passes directly through the construction. The transition maps and the limiting curve depend only on envelopes and compactness. A volume requires more: for a nonclosed background, the limiting Monge--Amp\`ere mass can change under a smooth change of representative.

Throughout this chapter, $X$ is a connected compact complex manifold of dimension $n$, $\theta$ is a smooth real $(1,1)$-form, and $a<c$ are fixed real numbers.

\section{Curves constant on a left half-line}
\label{sec:left-constant-curves}

Constancy on the left fixes the initial potential of the Legendre ray and bounds its growth between the two slopes $a$ and $c$. We require each slice to be fixed by the multiplier-ideal envelope of \cref{def:hermitian-i-envelope}.

\begin{definition}\label{def:bounded-left-i-curve}
Let $\mathscr T_{a,c}(\theta)$ be the set of families $\Gamma=(\Gamma_\tau)_{\tau<c}$ satisfying:
\begin{enumerate}
\item $\Gamma_\tau\in\PSH(X,\theta)$ and $P_\theta[\Gamma_\tau]_{\mathcal I}=\Gamma_\tau$;
\item the family is pointwise concave and decreasing in $\tau$;
\item $\Gamma_\tau=\Gamma_a$ for $\tau\leq a$.
\end{enumerate}
The cutoff is $c$, with the extension $\Gamma_c=\lim_{\tau\uparrow c}\Gamma_\tau$ and $\Gamma_\tau=-\infty$ for $\tau>c$. The order on these curves is pointwise order of their finite slices.
\end{definition}

Every slice has supremum zero. The decreasing endpoint limit is genuine by \cref{thm:normalized-compactness,cor:decreasing-limit}, and it is also $P_\theta[\cdot]_{\mathcal I}$-fixed: order preservation gives
\[
 \Gamma_c\leq P_\theta[\Gamma_c]_{\mathcal I}\leq\Gamma_\tau
 \quad(\tau<c),
\]
so taking the infimum proves equality. The first inequality uses the normalized potential itself as a competitor.

\begin{proposition}\label{prop:bounded-left-i-ray}
A curve in $\mathscr T_{a,c}(\theta)$ is an envelope test curve with initial potential $\varphi=\Gamma_a$. Denote its Legendre transform by $U_t=(\sup_{\tau<c}(\Gamma_\tau+t\tau))^*$ for real $t>0$. It is a Perron envelope ray and satisfies
\[
 \varphi+at\leq U_t\leq\varphi+ct\qquad(t>0).
\]
In particular $U_t\sim\varphi$ at each finite time.
\end{proposition}
\begin{proof}
A $P_\theta[\cdot]_{\mathcal I}$-fixed potential is $P_\theta$-fixed by \cref{cor:ordinary-preserves-i}. The regularized supremum of the curve is exactly $\Gamma_a$, so all the hypotheses of \cref{def:hermitian-envelope-testcurve} hold. Apply \cref{thm:testcurve-to-envelope-ray}. The term $\tau=a$ in the transform gives the lower bound, and $\Gamma_\tau\leq\varphi$, $\tau<c$, give the upper bound.
\end{proof}

\section{Transition maps}
\label{sec:envelope-system-transitions}

Projection to a larger background preserves every slice's multiplier ideals. Uniqueness of the normalized fixed representative then gives both injectivity and compatibility of successive projections.

\begin{proposition}\label{prop:hermitian-i-curve-transition}
If $\beta\geq\theta$ are smooth real backgrounds, slicewise projection defines an injective map
\[
 P_\beta[\cdot]_{\mathcal I}:\mathscr T_{a,c}(\theta)\longrightarrow
                       \mathscr T_{a,c}(\beta).
\]
For $\beta\geq\beta'\geq\theta$, these maps satisfy
\[
 P_\beta\bigl[P_{\beta'}[\Gamma]_{\mathcal I}\bigr]_{\mathcal I}=P_\beta[\Gamma]_{\mathcal I}.
\]
They preserve order.
\end{proposition}
\begin{proof}
The input slices are $\beta$-plurisubharmonic, so their projections are genuine and $P_\beta[\cdot]_{\mathcal I}$-fixed. Decreasingness, constancy for $\tau\leq a$, and order preservation follow from \cref{prop:i-envelope-operations}. To check concavity, choose real $r,t<c$ and $0<s<1$, and put $\tau=(1-s)r+s t$. That same proposition and the concavity of the original curve give
\[
 \begin{split}
 (1-s)P_\beta[\Gamma_r]_{\mathcal I}+sP_\beta[\Gamma_t]_{\mathcal I}
 &\leq P_\beta[(1-s)\Gamma_r+s\Gamma_t]_{\mathcal I}\\
 &\leq P_\beta[\Gamma_\tau]_{\mathcal I}.
 \end{split}
\]
Here the first inequality follows by combining the scaling and addition formulas with backgrounds $(1-s)\beta$ and $s\beta$.

Each projection preserves the multiplier ideals at every weight. Thus the two sides of the transition identity are $P_\beta[\cdot]_{\mathcal I}$-fixed representatives of the same multiplier-ideal class, and are equal. If two projected curves agree, their original slices have the same multiplier ideals; since those original slices are $P_\theta[\cdot]_{\mathcal I}$-fixed, they agree as well. This proves injectivity.
\end{proof}

\Needspace{18\baselineskip}
Let $\mathscr H(X)$ denote the set of smooth positive Hermitian forms on $X$.

\begin{definition}\label{def:hermitian-envelope-system}
An \emph{envelope system over $\theta$}, with parameters $a<c$, is a family
\[
 \boldsymbol\Gamma=(\Gamma^\eta)_{\eta\in\mathscr H(X)},
 \qquad \Gamma^\eta\in\mathscr T_{a,c}(\theta+\eta),
\]
satisfying
\begin{equation}\label{eq:hermitian-system-compatibility}
 \Gamma^\eta=P_{\theta+\eta}[\Gamma^{\eta'}]_{\mathcal I}
 \quad\hbox{whenever }\eta\geq\eta'.
\end{equation}
Write $\mathscr S_{a,c}(\theta)$ for this set, ordered componentwise. A system is \emph{realized in the background $\theta$} if it has the form
\[
 \mathsf E(\Gamma)^\eta=P_{\theta+\eta}[\Gamma]_{\mathcal I}
 \quad\hbox{for some }\Gamma\in\mathscr T_{a,c}(\theta).
\]
\end{definition}

The transition maps go from a smaller perturbation to a larger one. All components have the same multiplier ideals at each $\tau<c$: compare two of them with the component indexed by $\eta+\eta'$.

If $u\in\PSH(X,\theta)$, the constant family $\Gamma_\tau=P_\theta[u]_{\mathcal I}$ for $\tau<c$, with the specified cutoff, lies in $\mathscr T_{a,c}(\theta)$. Its projections form an envelope system. The transition proposition also proves that $\mathsf E$ is injective.

\section{A canonical limit in the original background}

Compatibility orders the slices as the background shrinks. Their normalization prevents collapse, and their decreasing limit is again fixed by the multiplier-ideal envelope. Projecting this limit back to the larger backgrounds gives the greatest realized system below the original one.

\Needspace{21\baselineskip}
\begin{theorem}\label{thm:hermitian-system-retraction}
Fix any Hermitian form $h$ on $X$. The index $j$ below ranges over the positive integers. For every $\boldsymbol\Gamma\in\mathscr S_{a,c}(\theta)$ the limits
\[
 (\mathsf R\boldsymbol\Gamma)_\tau
 =\lim_{j\to\infty}\Gamma^{h/j}_\tau\qquad(\tau<c)
\]
exist as decreasing genuine potentials and define an element of $\mathscr T_{a,c}(\theta)$. They are independent of $h$. The maps satisfy
\[
 \mathsf R\mathsf E=\mathrm{id},
 \qquad \mathsf E\mathsf R(\boldsymbol\Gamma)
                       \leq\boldsymbol\Gamma.
\]
Both maps preserve order. More precisely,
\begin{equation}\label{eq:hermitian-system-order-adjunction}
 \mathsf E(\Gamma)\leq\boldsymbol\Gamma
 \quad\Longleftrightarrow\quad
 \Gamma\leq\mathsf R\boldsymbol\Gamma.
\end{equation}
Thus $\mathsf E\mathsf R(\boldsymbol\Gamma)$ is the greatest system below $\boldsymbol\Gamma$ that is realized in the original background.
\end{theorem}
\begin{proof}
Fix $\tau<c$ and put $q_j=\Gamma^{h/j}_\tau$. Compatibility and the fact that $q_{j+1}\leq0$ is a competitor for its larger-background envelope give $q_j\geq q_{j+1}$. The functions have supremum zero and a common smooth lower curvature bound. By \cref{thm:normalized-compactness,cor:decreasing-limit}, their decreasing limit $q$ is genuine and has supremum zero. Passing the distributional inequalities
\[
 \theta+h/j+\ddc q_j\geq0
\]
to the limit gives $q\in\PSH(X,\theta)$.

We claim $P_\theta[q]_{\mathcal I}=q$. Since $q\leq q_j$, its multiplier ideals are contained in those of $q_j$. The potential $P_\theta[q]_{\mathcal I}$ belongs to the larger background $\theta+h/j$ and has the ideals of $q$. Hence order preservation there gives
\[
 P_\theta[q]_{\mathcal I}
 \leq P_{\theta+h/j}[P_\theta[q]_{\mathcal I}]_{\mathcal I}
 \leq P_{\theta+h/j}[q_j]_{\mathcal I}=q_j.
\]
Take the infimum in $j$ and use $P_\theta[q]_{\mathcal I}\geq q$ to prove the claim. Concavity, decreasingness and constancy on the left half-line pass pointwise to these limits. They therefore form a curve in $\mathscr T_{a,c}(\theta)$.

For every $\eta>0$ there is $j$ with $h/j\leq\eta$. Compatibility gives $\Gamma^\eta_\tau\geq\Gamma^{h/j}_\tau$. Consequently
\[
 (\mathsf R\boldsymbol\Gamma)_\tau
 =\inf_{\eta\in\mathscr H(X)}\Gamma^\eta_\tau,
\]
which proves independence of $h$. It also proves order preservation of $\mathsf R$.

Write $q=(\mathsf R\boldsymbol\Gamma)_\tau$. Since $q\leq\Gamma^\eta_\tau$, order preservation gives
$P_{\theta+\eta}[q]_{\mathcal I}\leq\Gamma^\eta_\tau$.
This proves $\mathsf E\mathsf R\leq\mathrm{id}$.

Now start with a base curve $\Gamma$. Its slice $\Gamma_\tau$ is a competitor in every larger background, so
\[
 \Gamma_\tau\leq(\mathsf R\mathsf E\Gamma)_\tau
                  \leq P_{\theta+\eta}[\Gamma_\tau]_{\mathcal I}
 \quad(\eta>0).
\]
The outside terms have the same multiplier ideals. By the two ideal inclusions supplied by these inequalities, the middle term has those ideals as well. Both it and $\Gamma_\tau$ are $P_\theta[\cdot]_{\mathcal I}$-fixed, so they are equal. This proves $\mathsf R\mathsf E=\mathrm{id}$.

Order preservation of $\mathsf E$ follows from the transition proposition. Applying $\mathsf R$ proves the forward implication in \eqref{eq:hermitian-system-order-adjunction}. For the reverse implication apply $\mathsf E$ and use $\mathsf E\mathsf R\leq\mathrm{id}$. This equivalence proves the final maximality assertion.
\end{proof}

\Needspace{13\baselineskip}
\begin{proposition}\label{prop:hermitian-system-realization-criterion}
A system $\boldsymbol\Gamma\in\mathscr S_{a,c}(\theta)$ is realized in $\theta$ if and only if, for every $\tau<c$, its limiting slice $(\mathsf R\boldsymbol\Gamma)_\tau$ is multiplier-ideal equivalent to one, and hence every, slice $\Gamma^\eta_\tau$.
\end{proposition}
\begin{proof}
For a realized system this follows from $\mathsf R\mathsf E=\mathrm{id}$ and ideal preservation by each projection. Conversely the assumed equivalence and fixedness in the larger background give
\[
 P_{\theta+\eta}[(\mathsf R\boldsymbol\Gamma)_\tau]_{\mathcal I}
                       =\Gamma^\eta_\tau.
\]
Thus the system equals $\mathsf E\mathsf R$ of itself.
\end{proof}

The criterion can be expressed entirely locally. By \cref{thm:local-i-inputs}, it asks that the Lelong numbers of the limiting slice and a component agree at every point on every local modification. The limiting slice lies below every component, so its Lelong numbers can only be larger. Equality is not automatic: \cref{sec:envelope-boundary-example} will exhibit an extra divisorial pole in the limit, even for a positive K\"ahler background.

\section{Smooth changes of representative}

Subtracting a smooth potential changes the background and leaves multiplier ideals unchanged. Projection restores the envelope normalization. The resulting bijection also respects the limiting curve, because the gauge estimate is uniform over all perturbations.

\begin{proposition}\label{prop:hermitian-system-gauge}
If $\theta'=\theta+\ddc f$ with $f\in C^\infty(X,\mathbb R)$, take $\Gamma\in\mathscr T_{a,c}(\theta)$ and $\boldsymbol\Gamma\in\mathscr S_{a,c}(\theta)$. The formulas
\[
 (\mathsf G_f\Gamma)_\tau=P_{\theta'}[\Gamma_\tau-f]_{\mathcal I},
 \qquad
 (\mathsf G_f\boldsymbol\Gamma)^\eta_\tau
                    =P_{\theta'+\eta}[\Gamma^\eta_\tau-f]_{\mathcal I}
\]
define bijections on the corresponding curve and system spaces. They satisfy the cocycle law and commute with $\mathsf E$ and $\mathsf R$.
\end{proposition}
\begin{proof}
Subtracting $f$ changes the background as indicated and preserves concavity. Slicewise projection preserves concavity by the same convex-combination argument as before. It preserves the cutoff and the constant left part. Every new slice has the multiplier ideals of the original one because $f$ is smooth. Applying the construction with $-f$ gives a $P_\theta[\cdot]_{\mathcal I}$-fixed slice with those same ideals, hence the original slice. This proves the inverse statement; the same uniqueness proves the cocycle law.

Compatibility of the system and commutation with $\mathsf E$ follow by comparing multiplier ideals in the common final background. For commutation with $\mathsf R$, the gauge estimate of \cref{prop:i-envelope-gauge}, applied in each background $\theta+\eta$, gives
\[
 \Gamma^\eta_\tau+\min_Xf
 \leq (\mathsf G_f\boldsymbol\Gamma)^\eta_\tau+f
 \leq\Gamma^\eta_\tau+\max_Xf.
\]
The constants are independent of $\eta$. Pass to the cofinal decreasing limits. The new limit is ordinarily equivalent to the old limit minus $f$, hence has the same multiplier ideals. It is $P_{\theta'}[\cdot]_{\mathcal I}$-fixed by \cref{thm:hermitian-system-retraction}. It must therefore equal $P_{\theta'}[(\mathsf R\boldsymbol\Gamma)_\tau-f]_{\mathcal I}$, proving commutation with $\mathsf R$.
\end{proof}

When $X$ is K\"ahler, the closed positive perturbations are cofinal among all Hermitian perturbations: a small positive multiple of one fixed K\"ahler form is below any given Hermitian form. Hence the compatibility formula determines an envelope system from its closed-positive components. Conversely, a compatible family on that cofinal collection extends uniquely: for $\eta>0$ choose a closed positive $\kappa\leq\eta$ and set $\Gamma^\eta=P_{\theta+\eta}[\Gamma^\kappa]_{\mathcal I}$. Two choices give the same result by passing to a smaller closed positive form and using the transition identity.

For a closed background on a compact K\"ahler manifold, this is the indexing used in \cite[chapter ``Non-Archimedean pluripotential theory'', section ``The definition of non-Archimedean metrics'']{SGPT}. That construction allows general $\mathcal I$-model test curves and requires a positive-mass initial potential in each perturbed background. Our curves are constant for $\tau\leq a$; no mass condition has entered the system space or its retraction.

\section{Mass under a smooth change of background}

The gauge bijection preserves the singularity data of every component. For closed backgrounds, the volume construction in \cite[section ``The definition of non-Archimedean metrics'']{SGPT} takes a limit of the masses of the initial potentials. For a nonclosed background, even the zero-slice system detects the mass variation of \cref{ex:torus-mass}.

\Needspace{13\baselineskip}
\begin{proposition}\label{prop:envelope-system-raw-volume-failure}
There are Hermitian forms $\omega,h$ on a compact complex surface and a smooth real function $f$, with $\omega'=\omega+\ddc f>0$, for which the following holds. For every $a<c$ there is a system $\boldsymbol\Gamma\in\mathscr S_{a,c}(\omega)$ such that both limits
\[
 \lim_{\epsilon\downarrow0}
 \int_X\mathrm{MA}^{\mathrm{ext}}_{\omega+\epsilon h}
          (\Gamma^{\epsilon h}_a),
 \qquad
 \lim_{\epsilon\downarrow0}
 \int_X\mathrm{MA}^{\mathrm{ext}}_{\omega'+\epsilon h}
          ((\mathsf G_f\boldsymbol\Gamma)^{\epsilon h}_a)
\]
exist and are unequal.
\end{proposition}
\begin{proof}
Take the torus and the smooth perturbation from \cref{prop:raw-i-volume-gauge-failure}, so that $\omega'=\omega+\ddc f$ is Hermitian and $\int_X(\omega')^2\neq\int_X\omega^2$. Fix any Hermitian form $h$. In every background $\omega+\eta$, let every slice with $\tau<c$ be zero. This is an envelope system for any $a<c$: zero is the multiplier-ideal envelope of the unit-ideal class in a positive background. Its first limit is
\[
 \lim_{\epsilon\downarrow0}\int_X(\omega+\epsilon h)^2
                          =\int_X\omega^2.
\]
The gauge bijection sends every slice to $P_{\omega'+\eta}[-f]_{\mathcal I}=0$, since $f$ is smooth and the new background is positive. The second limit is therefore $\int_X(\omega')^2$, which is different.
\end{proof}

Thus this limit of raw masses does not give a volume invariant under smooth changes of representative. Recovery of the system poses a separate question, already on a K\"ahler surface, where closed-background masses are invariant.

\section{Failure of recovery in a positive background}
\label{sec:envelope-boundary-example}

The realization condition in \cref{prop:hermitian-system-realization-criterion} is not automatic, even when the original background is K\"ahler. The obstruction can occur because the singularity leaves a degenerate residual class. We give a direct reconstruction of the ruled-surface example of Demailly--Peternell--Schneider \cite[Example 1.7]{DPS94}, then apply it to envelope systems.

In this section the general manifold is replaced by the following surface. Let $E=\C/(\mathbb Z+\mathrm i\mathbb Z)$ and form the compact ruled surface
\[
 X=(\C_z\times\mathbb P^1_w)/\sim,
 \qquad (z,w)\sim(z+1,w),\qquad
 (z,w)\sim(z+\mathrm i,w+1).
\]
The transformations commute and cover the deck transformations of $E$. They define a holomorphic $\mathbb P^1$-bundle $\pi:X\to E$. Let $D$ be its section $w=\infty$. On $X\setminus D$, the smooth function $v=w-\operatorname{Im}z$ is invariant under both transformations.

\Needspace{12\baselineskip}
\begin{lemma}\label{lem:ruled-explicit-nef}
Let $\sigma_0=(\mathrm i/2)\,\mathrm{d}z\wedge\mathrm{d}\bar z$ be the flat probability form on $E$. There is a smooth closed real $(1,1)$-form $\theta\in c_1(\mathcal O_X(D))$ whose class is nef, and a K\"ahler form
\[
 \kappa=\theta+\pi^*\sigma_0,
 \qquad \sigma_0=\frac{\mathrm i}{2}\,\mathrm{d}z\wedge \mathrm{d}\bar z,
 \qquad \int_E\sigma_0=1.
\]
The only positive current of the form $\theta+\ddc u$ is $[D]$.
\end{lemma}
\begin{proof}
Let $s_D$ be the canonical section of $\mathcal O_X(D)$ with zero divisor $D$. For a real parameter $t>0$, define a smooth Hermitian metric on this line bundle by prescribing its norm:
\[
 |s_D|_t^2=(1+t|v|^2)^{-1}\quad\hbox{on }X\setminus D.
\]
To check smooth extension, use $\xi=1/w$ near $D$. The local squared norm of a nonvanishing frame is
\[
 \bigl(|\xi|^2+t|1-(\operatorname{Im}z)\xi|^2\bigr)^{-1},
\]
which is smooth and positive. The displayed norm is invariant on the complement, so these extensions define a global metric. Its curvature is
\[
 \theta_t=\ddc\log(1+t|v|^2)
 \quad\hbox{on }X\setminus D.
\]
Set $\theta=\theta_1$ and $\sigma=\sigma_0/\pi$ on $E$, where the denominator is the number $\pi$.

Here is an explicit curvature estimate. Put $r=|v|^2$ and let $H$ be the complex Hessian of $\log(1+tr)$ in $(z,w)$. Direct differentiation gives
\[
 \begin{split}
 H_{w\bar w}&=\frac{t}{(1+tr)^2},\\
 H_{z\bar w}&=\frac{\mathrm i t(1-tv^2)}{2(1+tr)^2},\\
 H_{z\bar z}&=\frac{t}{2(1+tr)}
             -\frac{t^2(\operatorname{Re}v)^2}{(1+tr)^2}.
 \end{split}
\]
Consequently
\[
 H_{z\bar z}-\frac{|H_{z\bar w}|^2}{H_{w\bar w}}
 =\frac{t(1-tr)}{4(1+tr)}\geq-\frac t4.
\]
Completing the square proves $\theta_t\geq-(t/4)\pi^*\sigma$. This inequality extends smoothly across $D$. The restriction of $\theta_t$ to every fiber is a positive Fubini--Study form, including at infinity. Thus $\theta_t+A\pi^*\sigma$ is strictly positive off $D$ whenever $A>t/4$. At $\xi=0$, its coefficient matrix, with the common factor $\mathrm i/(2\pi)$ omitted, is
\[
 \begin{pmatrix} A&\mathrm i/2\\-\mathrm i/2&1/t\end{pmatrix}.
\]
Its vertical coefficient is positive and its Schur complement is $A-t/4>0$. This checks strict positivity also at infinity. In particular $\kappa=\theta+\pi^*\sigma_0$ is K\"ahler.

\Needspace{7\baselineskip}
The function
\[
 f_t=\log\frac{1+t|v|^2}{1+|v|^2}
\]
extends smoothly across $D$ and satisfies $\theta_t=\theta+\ddc f_t$. The curvature lower bound, with $t\downarrow0$, proves nefness in the smooth metric definition.

For the rigidity assertion we use a one-variable observation. For real constants $A,B$, a genuine entire subharmonic function $h$ on $\C$, with coordinate $\zeta$, satisfying $h(\zeta+1)=h(\zeta)$ and
\[
 h(\zeta)\leq A\log(1+|\zeta|)+B
\]
is constant. Indeed the upper-semicontinuous regularization of
$\sup_{s\in[0,1]}h(\zeta+s)$ is subharmonic and invariant under every real translation. It is therefore a convex function of $\operatorname{Im}\zeta$. Its upper growth is sublinear at both ends, so convexity makes it constant. Hence $h$ is bounded above, and the subharmonic Liouville theorem makes $h$ constant as well.

Write $\psi_D=\log|s_D|_1^2$. Poincar\'e--Lelong gives $\theta+\ddc\psi_D=[D]$. If $u\in\PSH(X,\theta)$, then $u-\psi_D$ is plurisubharmonic on $X\setminus D$. Its lift $H(z,w)$ to $\C^2$ is invariant under the two deck transformations and, since $u$ is bounded above on compact $X$, satisfies the following inequality for some real constant $C$:
\[
 H(z,w)\leq C+\log\bigl(1+|w-\operatorname{Im}z|^2\bigr).
\]
For each fixed $w$, the restriction in $z$ is either identically $-\infty$ or a genuine subharmonic function. It is $1$-periodic and has logarithmic upper growth, so the observation makes it constant in $z$. Thus $H(z,w)=h(w)$, with $h$ a genuine subharmonic function of $w$. The second deck transformation now gives $h(w+1)=h(w)$, and the bound at $z=0$ gives logarithmic growth. The same observation makes $h$ constant. Therefore $u=\psi_D+C$ off $D$, and equality of the canonical quasi-plurisubharmonic representatives extends it across $D$. This proves the assertion.
\end{proof}

Fix a point $p\in E$ and let $F=\pi^{-1}(p)$. The Green-potential construction on curves gives a quasi-plurisubharmonic $g$ with
$\sigma_0+\ddc g=\delta_p$. Its pullback $\psi_F=g\circ\pi$ has a logarithmic pole of coefficient one along $F$, and
\[
 \pi^*\sigma_0+\ddc\psi_F=[F].
\]

\Needspace{20\baselineskip}
\begin{theorem}\label{thm:positive-background-system-nonrecovery}
For the K\"ahler background $\kappa$ above and every $a<c$, there is a system $\boldsymbol\Gamma\in\mathscr S_{a,c}(\kappa)$ for which
\[
 \mathsf E\mathsf R(\boldsymbol\Gamma)
                         \neq\boldsymbol\Gamma.
\]
Every finite slice of every component has multiplier ideals
\[
 \mathcal I(t\Gamma^\eta_\tau)
                 =\mathcal O_X(-\lfloor t\rfloor F)\qquad(t>0),
\]
whereas its limiting slice has ideals
\[
 \mathcal I\bigl(t(\mathsf R\boldsymbol\Gamma)_\tau\bigr)
                 =\mathcal O_X(-\lfloor t\rfloor(F+D)).
\]
In particular $\mathsf E$ is not surjective and $\mathsf R$ is not injective on these unrestricted system spaces.
\end{theorem}
\begin{proof}
For each Hermitian form $\eta>0$, the explicit nef approximation supplies a smooth $f_\eta$ such that
$\theta+\eta+\ddc f_\eta\geq0$. Then
\[
 u_\eta=\psi_F+f_\eta\in\PSH(X,\kappa+\eta).
\]
Define every finite slice of $\Gamma^\eta$ to be
$q_\eta=P_{\kappa+\eta}[u_\eta]_{\mathcal I}$, with cutoff $c$. Its multiplier ideals are those of $\psi_F$, independent of the smooth choice $f_\eta$. Hence the transition identity follows by comparing these ideals in the larger background. The constant curves form a system with the required parameters.

Let $q$ be its limiting slice under $\mathsf R$, and let $\nu(q,F)$ denote its generic Lelong number along $F$. Each $q_\eta$ has generic Lelong number one along $F$, and upper semicontinuity under the decreasing cofinal limit gives $\nu(q,F)\geq1$. Divisorial removal therefore gives
\[
 \kappa+\ddc q-[F]\geq0.
\]
By Poincar\'e--Lelong this current is $\theta+\ddc(q-\psi_F)$, where the difference is first taken as an integrable distribution. Positivity supplies its canonical $\theta$-plurisubharmonic representative. The rigidity lemma makes this current $[D]$, so
\[
 q=\psi_F+\psi_D+C.
\]
The two divisors are smooth and meet transversely, since one is a fiber and the other a section. Direct product integrability in local coordinates gives the stated ideals for every real $t>0$. At a generic point of $D$, the Lelong number has thus increased from zero to one. The realization criterion proves the strict failure of recovery.

Finally, $\mathsf R\mathsf E=\mathrm{id}$ implies that $\boldsymbol\Gamma$ and the distinct realized system $\mathsf E\mathsf R\boldsymbol\Gamma$ have the same image under $\mathsf R$. This proves its failure of injectivity.
\end{proof}

\Needspace{25\baselineskip}
\begin{proposition}\label{prop:nonrecovery-zero-residual-volume}
Let $q_\eta$ be the common finite slice constructed in the preceding proof. Along the closed perturbations $\eta=\epsilon\kappa$, with real $\epsilon>0$, the component slices of this example satisfy
\[
 \int_X\mathrm{MA}^{\mathrm{ext}}_{\kappa+\epsilon\kappa}(q_{\epsilon\kappa})
 =2\epsilon\int_X\theta\wedge\kappa
                  +\epsilon^2\int_X\kappa^2\longrightarrow0.
\]
The limiting slice has zero non-pluripolar Monge--Amp\`ere measure.
\end{proposition}
\begin{proof}
The defining competitor gives
$q_\eta\geq\psi_F+f_\eta-\sup_X(\psi_F+f_\eta)$.
On the other hand, removing its divisorial component $[F]$ gives a positive current in background $\theta+\eta$, with canonical global quasi-plurisubharmonic potential $w_\eta=q_\eta-\psi_F$. This potential is bounded below by the displayed inequality and bounded above by compactness. Hence $q_\eta\simeq\psi_F$ and the residual potential $w_\eta$ is bounded.

Off $F$, plurifine locality identifies the Monge--Amp\`ere measure with $(\theta+\eta+\ddc w_\eta)^2$. For closed $\eta=\epsilon\kappa$ the latter bounded product assigns zero mass to $F$ and has total mass $\int_X(\theta+\epsilon\kappa)^2$, by bounded Stokes in a closed background.

The normal bundle of $D$ is trivial. Indeed the transition $w\mapsto w+1$ gives $\xi\mapsto\xi/(1+\xi)$, whose derivative at $\xi=0$ is one; the other transition is the identity. Thus
$\int_X\theta^2=D^2=0$. Expanding the square gives the formula. The limit potential has current $[F]+[D]$, which vanishes away from $F\cup D$. Its non-pluripolar measure is consequently zero.
\end{proof}

This example concerns a positive original background on a K\"ahler surface; failure of the bounded mass property is not its cause. The residual volume of the prescribed singularity tends to zero. The full inverse-system space in \cite{SGPT} permits zero volume; its recovery theorem from positive-volume base curves imposes a positive-volume hypothesis. The example is consistent with that distinction. A recovery theorem with a suitable positive-volume condition remains a different question in the Hermitian setting.

